% Extracted from p08-distillation
\subsection{Distilling into Recurrent Students}
\label{ch-distill}

\paragraph{Background}

\paragraph{Knowledge distillation.} \citet{hinton2015distilling}
introduced training a student on the teacher's softened output
distribution, transferring ``dark knowledge'' about class similarities.
\citet{sanh2019distilbert} extended this with hidden-state matching
between aligned layers, producing a widely used compact BERT.

\paragraph{Recurrent transformers.} The Universal Transformer
\citep{dehghani2019universal} replaces fixed depth with a shared block
applied recurrently. The OpenMythos variant frames the recurrent core
with unshared prelude layers (input processing) and coda layers
(output processing). Our reference student configuration, matching the
44M model studied in a companion paper of this series, uses dimension
512, 6 prelude plus 2 coda layers, 4 loops of a shared block, 8 experts
per MoE layer, grouped-query attention, and a 50,257-token BPE
vocabulary (43,974,722 parameters total), trained on TinyStories
\citep{eldan2023tinystories}.

\paragraph{The calibration problem.} In the companion study, the 44M
model's test-time loop sweep showed perplexity 29.38 at 1 loop falling
to 4.41 at the trained depth of 4 loops, then degrading slightly to
4.90 at 8 loops and plateauing through 16 loops. The recurrent dynamics
are calibrated to the training iteration count; extra test-time loops
do not refine further. Any distillation protocol for recurrent students
must reckon with this: supervision that fixes a single loop count
risks baking in the same calibration.

\paragraph{Why Naive Distillation Under-Constrains Recurrence}

Consider a teacher $T$ with $L$ distinct layers and a student $S$ whose
shared block $f_\theta$ is iterated $K$ times between prelude $P$ and
coda $C$:
\begin{equation}
h_0 = P(x), \quad h_{k} = f_\theta(h_{k-1})\;\; (k=1\ldots K), \quad
y = C(h_K).
\end{equation}
Naive distillation minimises a divergence $D(p_T(y\mid x)\,\|\,p_S(y\mid
x))$ over the final outputs only. The difficulty is degrees of
freedom: for fixed $P$, $C$, and output $y$, the intermediate sequence
$(h_1, \ldots, h_{K-1})$ is unconstrained. In particular, the
degenerate solution $f_\theta \approx \mathrm{identity}$ satisfies any
output-matching objective while performing zero iterative refinement;
all computation migrates to $P$ and $C$, which the optimiser prefers
because they receive undiluted gradients. The student then matches the
teacher's logits with a recurrent block that does nothing---and, per
the calibration finding above, its loop count becomes a frozen
hyperparameter rather than a meaningful compute knob.

A second, subtler failure: even when $f_\theta$ is non-trivial, nothing
in the objective requires successive iterations to implement
\emph{progressive} refinement. The block may learn oscillatory or
redundant dynamics that happen to land near the teacher's output at
exactly step $K$. Attention-map inspection in such students (proposed
in \S\ref{sec:experiments}) would be expected to show near-identical
patterns across loops, in contrast to a teacher whose layer-wise
attention evolves from local to global structure.

\paragraph{Loop-Aware Distillation Protocol}

We propose supervising the \emph{trajectory}, not just the endpoint.

\paragraph{Layer-to-loop alignment.} Partition the teacher's $L$
layers into $K$ contiguous groups $G_1, \ldots, G_K$. For loop $k$,
define the target representation as the teacher's output after group
$G_k$:
\begin{equation}
\mathcal{L}_{\mathrm{hid}} = \sum_{k=1}^{K} \lambda_k \,
\big\lVert h_k - \phi_k\big(T_{G_k}(x)\big) \big\rVert^2,
\end{equation}
where $T_{G_k}(x)$ is the teacher hidden state after group $k$ and
$\phi_k$ is a learned linear projection bridging dimensionality gaps.
The weights $\lambda_k$ follow a ramp $\lambda_k \propto k/K$,
encoding the prior that later iterations should carry more of the
teacher's high-level computation.

\paragraph{Cross-iteration attention transfer.} For each loop $k$ and
each attention head, we match the student's attention distribution to
the mean attention map of teacher group $G_k$ via KL divergence:
\begin{equation}
\mathcal{L}_{\mathrm{attn}} = \sum_{k=1}^{K} \mu_k \,
\mathrm{KL}\big(\bar{A}^{T}_{G_k} \,\|\, A^{S}_{k}\big).
\end{equation}
This directly pressures the recurrent block to reproduce the
teacher's layer-wise attention evolution---from local patterns early
to long-range structure late---rather than collapsing to static maps.

\paragraph{Loop-count curriculum.} To avoid calibrating the dynamics to
a single $K$, training randomises the executed loop count per batch,
$K \sim \mathcal{U}\{1, K_{\max}\}$, with the alignment targets
recomputed for the sampled depth (teacher groups re-partitioned
accordingly). Early training uses small $K_{\max}$, growing to the
full depth---a curriculum over iteration count. Final-logit matching
against the teacher is retained throughout:
\begin{equation}
\mathcal{L} = \alpha \mathcal{L}_{\mathrm{task}}
+ \beta \mathcal{L}_{\mathrm{logit}}
+ \gamma \mathcal{L}_{\mathrm{hid}}
+ \delta \mathcal{L}_{\mathrm{attn}},
\end{equation}
with $\gamma, \delta$ warmed up from zero once the student matches the
teacher's outputs coarsely (we propose $\alpha{=}1, \beta{=}1$,
$\gamma, \delta$ ramping to $0.5$).

\paragraph{Proposed Experiments}
\label{sec:experiments}

\paragraph{Setup.} Teacher: a standard 12-layer transformer
($\approx$120M parameters, proposed configuration) trained on
TinyStories \citep{eldan2023tinystories}. Student: the 44M recurrent
configuration above ($K_{\max}{=}4$). All runs proposed; no results
are claimed.

\paragraph{Experiment 1: ablation of supervision terms.} Train four
students: (a) naive ($\beta$ only); (b) +$\mathcal{L}_{\mathrm{hid}}$;
(c) +$\mathcal{L}_{\mathrm{attn}}$; (d) full protocol. Report
perplexity at each loop count $1\ldots8$ and the loop-sweep curve
shape. \emph{Predicted} outcome: (a) matches teacher perplexity at
exactly 4 loops but shows flat or degraded behaviour off the trained
depth with static cross-loop attention; (d) shows monotonic
improvement with loop count and evolving attention maps. These are
hypotheses to be tested, not findings.

\paragraph{Experiment 2: iteration-content probes.} For each student,
train linear probes on $h_k$ to predict the next token, and measure
per-loop probe accuracy. A student doing genuine iterative refinement
should show rising probe accuracy across $k$; a degenerate student
will show flat probe curves with all gains in the coda.

\paragraph{Experiment 3: loop-count generalisation.} Evaluate each
student at $K \in \{1, 2, 4, 8, 12, 16\}$. The curriculum-trained
student is hypothesised to degrade gracefully rather than exhibiting
the sharp calibration seen in fixed-$K$ training.

\paragraph{Discussion}

If the protocol works, it yields a genuinely new capability: small
recurrent students that inherit a large teacher's knowledge
\emph{and} implement it as iterative refinement, with loop count as a
working test-time knob. If per-loop supervision proves unnecessary---if
naive distillation already produces refining dynamics---that is itself
an informative negative result about how much structure output
matching imposes.

A limitation is teacher choice: distilling from a fixed-depth teacher
may transfer depth-shaped computation that recurrence can only
approximate. An alternative worth testing is a recurrent teacher,
which would make the alignment exact but sacrifices the abundance of
strong pretrained fixed-depth teachers.

\paragraph{Limitations}

This is a method proposal without empirical validation. The predicted
outcomes in \S\ref{sec:experiments} are hypotheses. Scale is limited
to the sub-100M regime and TinyStories; whether loop-aware
distillation helps at larger scales or on reasoning-heavy data is
open. The alignment between teacher layer groups and student loops is
a heuristic---there may be no faithful mapping, in which case the
per-loop targets could harm more than help.

\paragraph{Conclusion}

Distilling into recurrent students requires supervising the
trajectory, not just the endpoint. We have specified a loop-aware
protocol---per-loop hidden targets, cross-iteration attention
transfer, and a loop-count curriculum---together with the experiments
that would test it. The question it answers is precise: can
distillation produce recurrent students whose loops do real work?
